\documentclass[11pt]{article}
\usepackage[T1]{fontenc}
\usepackage{lmodern}
\usepackage{amsmath,amssymb,amsthm,mathrsfs}
\usepackage{graphicx}
\usepackage{booktabs}
\usepackage{microtype}
\usepackage[margin=0.92in]{geometry}
\usepackage{xcolor}
\usepackage[colorlinks=true,linkcolor=blue!55!black,citecolor=blue!55!black,urlcolor=blue!55!black]{hyperref}

\newtheorem{theorem}{Theorem}
\newtheorem{lemma}{Lemma}
\newtheorem{corollary}{Corollary}
\theoremstyle{remark}
\newtheorem{remark}{Remark}
\newcommand{\Bd}{\mathbb B^d}
\newcommand{\Sd}{\partial\mathbb B^d}
\newcommand{\Eeta}{\mathcal E_{\eta}}
\newcommand{\Ieta}{\mathcal I_{\eta}}
\newcommand{\Tr}{\operatorname{Tr}}

\title{Arbitrary Poisson Boundary Measures Give Coordinate Multipliers}
\author{Candidate full-result packet for arXiv:2001.09616}
\date{26 August 2026}

\begin{document}
\maketitle

\begin{center}
\fbox{\parbox{0.91\textwidth}{
\textbf{Status: full solution likely valid; human review required.}
For every finite positive Borel measure on the sphere, the associated
Poisson-weighted volume measure is Carleson for the Dirichlet-type space in
the source.  Hence every coordinate function is a multiplier, and the
multiplication tuple is a joint $2$-isometry.  The proof is analytic and does
not assume polynomial density.  Novelty confidence is moderate after a
bounded search found no explicit later answer.
}}
\end{center}

\section{Source target and exact scope}

The source is S.~Chavan, R.~Gupta, and Md.~R.~Reza,
\emph{Dirichlet-type spaces on the unit ball and joint $2$-isometries},
arXiv:2001.09616v2 and J.~Funct.~Anal.~279 (2020), 108733
\cite{CGR2020}.  Let $d\geq2$, let $V$ be normalized volume measure on
$\Bd$, and let $\sigma$ be normalized surface measure on $\Sd$.  The source
uses the Euclidean Poisson kernel for the real $2d$-ball,
\begin{equation}\label{eq:poisson}
 P(z,\eta)=\frac{1-\lVert z\rVert^2}{\lVert z-\eta\rVert^{2d}},
 \qquad z\in\Bd,\quad \eta\in\Sd,
\end{equation}
and $P[\mu](z)=\int_{\Sd}P(z,\eta)\,d\mu(\eta)$.

For a finite positive Borel measure $\mu$ on $\Sd$, its scalar
Dirichlet-type space consists of the holomorphic functions with
\begin{equation}\label{eq:Dnorm}
 \lVert f\rVert_{\mathscr D(\mu)}^2
 =\lVert f\rVert_{H^2(\Bd)}^2
 +\frac1d\int_{\Bd}\lVert\nabla f(z)\rVert^2P[\mu](z)\,dV(z)<\infty.
\end{equation}
Published Proposition 3.5 says that all coordinate functions are multipliers
if and only if
\begin{equation}\label{eq:targetembedding}
 \int_{\Bd}|f(z)|^2P[\mu](z)\,dV(z)
 \leq C_\mu\lVert f\rVert_{\mathscr D(\mu)}^2.
\end{equation}
Page 12 proves this only when $\mu$ has bounded density and then states:
\begin{quote}
``We do not know whether the weighted volume measure with weight being the
Poisson integral of an arbitrary finite positive Borel measure $\mu$ on
$\partial\mathbb B^d$, $d\geq2$ is a Carleson measure for
$\mathscr D(\mu)$.''
\end{quote}
The theorem below proves \eqref{eq:targetembedding} for every such $\mu$.

\begin{figure}[ht]
\centering
\includegraphics[width=0.97\textwidth]{figures/open_problem_crop.png}
\caption{Source p.~12: the bounded-density corollary and the unresolved
arbitrary-measure statement.}
\end{figure}

\section{Proof intuition}

The Poisson integral is linear in the boundary measure, so the decisive case
is one atom $\delta_\eta$.  Its density looks highly singular in the usual
coordinates.  Center real polar coordinates at the singular boundary point
$\eta$.  The singular power cancels the polar Jacobian exactly:
\[
 P(z,\eta)\,dV(z)
 =c_{2d}(2\langle\omega,\eta\rangle_{\mathbb R}-\rho)
 \,d\rho\,d\omega,
 \qquad z=\eta-\rho\omega.
\]
Thus the measure has effective radial dimension one at its tip.  A weighted
one-dimensional Hardy estimate controls the atomic weighted $L^2$ norm by
the tip trace and the atomic gradient energy.

It remains to bound that trace.  Average $f$ over the unitary rotations that
fix $\eta$.  This preserves the radial value and contracts both the Hardy norm
and the atomic energy, while reducing $f$ to a one-variable function
$h(\langle z,\eta\rangle)$.  After the transverse variables are integrated
out, the energy density of $h'$ is at least $c_d/(1-r)$ on every dyadic
nontangential box at $1$.  Submean estimates and a bounded-overlap dyadic sum
then bound the trace.  Integrating the resulting uniform atomic inequality
against $\mu$ proves the full statement.

\section{Atomic geometry and trace estimates}

For $\eta\in\Sd$ define, with extended values allowed,
\begin{align}
 \Eeta(f)&=\int_{\Bd}\lVert\nabla f(z)\rVert^2P(z,\eta)\,dV(z),\label{eq:Eeta}\\
 \Ieta(f)&=\int_{\Bd}|f(z)|^2P(z,\eta)\,dV(z).\label{eq:Ieta}
\end{align}

\begin{lemma}[Tip trace and anchored atomic Poincare inequality]\label{lem:tip}
Let $f$ be holomorphic on $\Bd$ and suppose $\Eeta(f)<\infty$.  Then $f$ has
a single finite-energy trace $\Tr_\eta f$ at the tip $\eta$: along almost
every inward ray from $\eta$ the limit exists and equals the same number.
Moreover,
\begin{equation}\label{eq:bulkfromtrace}
 \Ieta(f)\leq 2|\Tr_\eta f|^2+\frac32\Eeta(f).
\end{equation}
\end{lemma}

\begin{proof}
Regard $\mathbb C^d$ as $\mathbb R^N$, $N=2d$, and put
$a=\langle\omega,\eta\rangle_{\mathbb R}$.  The coordinates
\[
 z=\eta-\rho\omega,\qquad a>0,\qquad 0<\rho<2a,
\]
parametrize the ball up to null sets.  Indeed,
\[
 1-|z|^2=\rho(2a-\rho),\qquad |z-\eta|=\rho.
\]
Since $dV=c_N\rho^{N-1}d\rho\,d\omega$, the exact cancellation is
\begin{equation}\label{eq:jacobian}
 P(z,\eta)dV(z)=c_N(2a-\rho)d\rho\,d\omega.
\end{equation}

We first justify the common trace.  On every compact connected angular set
$K\Subset\{a>0\}$ and for all sufficiently small $\rho$, the density in
\eqref{eq:jacobian} is bounded below.  The radial part of the energy implies
that $f(\eta-\rho\omega)$ has a limit $L(\omega)$ for almost every
$\omega\in K$, and convergence holds in $L^2(K)$.  In centered polar
coordinates the angular part of the Euclidean gradient is
$\rho^{-1}\nabla_\omega$.  Hence
\[
 \int_0^{\rho_0}\rho^{-2}
 \int_K|\nabla_\omega f(\eta-\rho\omega)|^2d\omega\,d\rho<\infty.
\]
There is a sequence $\rho_n\downarrow0$ on which the inner angular energy
tends to zero.  The Poincare inequality on $K$, followed by the $L^2(K)$
convergence, shows that $L$ is constant on $K$.  Exhausting the connected
hemisphere $\{a>0\}$ by overlapping compact sets gives one common value,
denoted $\Tr_\eta f$.

Fix a direction and write $L=2a$ and
$u(\rho)=f(\eta-\rho\omega)$.  The fundamental theorem of calculus and
Cauchy--Schwarz give
\[
 |u(\rho)-\Tr_\eta f|^2
 \leq \rho\int_0^\rho|u'(s)|^2ds.
\]
After multiplication by $L-\rho$ and integration, the coefficient of
$|u'(s)|^2$ is
\[
 A_L(s)=\int_s^L\rho(L-\rho)d\rho.
\]
Writing $t=s/L$, direct factorization gives
\[
 \frac{A_L(s)}{L^2(L-s)}
 =\frac{(1-t)(2t+1)}6\leq\frac3{16}.
\]
Because $L\leq2$,
\begin{equation}\label{eq:rayhardy}
 \int_0^L(L-\rho)|u(\rho)-\Tr_\eta f|^2d\rho
 \leq\frac34\int_0^L(L-\rho)|u'(\rho)|^2d\rho.
\end{equation}
Integrate over $\omega$, use \eqref{eq:jacobian}, and note that
$|u'|\leq\lVert\nabla f\rVert$.  Since
$\int_{\Bd}P(z,\eta)dV(z)=1$, the inequality
$|u|^2\leq2|\Tr_\eta f|^2+2|u-\Tr_\eta f|^2$ and
\eqref{eq:rayhardy} yield \eqref{eq:bulkfromtrace}.
\end{proof}

\begin{lemma}[Uniform holomorphic tip trace]\label{lem:trace}
For every $d\geq2$ there is $C_d<\infty$ such that, for every
$\eta\in\Sd$ and every holomorphic $f$ with
$f\in H^2(\Bd)$ and $\Eeta(f)<\infty$,
\begin{equation}\label{eq:tracebound}
 |\Tr_\eta f|^2\leq C_d\bigl(\lVert f\rVert_{H^2(\Bd)}^2+\Eeta(f)\bigr).
\end{equation}
\end{lemma}

\begin{proof}
Unitary invariance reduces to $\eta=e_1$.  Let $G\simeq U(d-1)$ be the
compact group of unitaries fixing $e_1$, and let
\[
 (Qf)(z)=\int_G f(Uz)\,dU.
\]
Haar averaging and Jensen's inequality show
\begin{equation}\label{eq:contractions}
 \lVert Qf\rVert_{H^2}\leq\lVert f\rVert_{H^2},
 \qquad \mathcal E_{e_1}(Qf)\leq\mathcal E_{e_1}(f).
\end{equation}
The holomorphic $G$-invariants are precisely the functions of the first
coordinate, so $Qf(z)=h(z_1)$, where $h(\lambda)=f(\lambda,0,\ldots,0)$.
The tip trace is invariant under $G$ and linear on the finite-energy class;
hence $\Tr_{e_1}Qf=\Tr_{e_1}f$.

Put $A=1-|\lambda|^2$ and $b=|1-\lambda|^2$.  Integrating the other
$d-1$ complex coordinates in polar coordinates gives, up to a positive
dimensional normalization,
\begin{equation}\label{eq:marginal}
 \mathcal E_{e_1}(Qf)
 =c_d\int_{\mathbb D}|h'(\lambda)|^2W_d(\lambda)dA(\lambda),
 \qquad
 W_d(\lambda)=\int_0^A
 \frac{(A-u)u^{d-2}}{(b+u)^d}\,du.
\end{equation}

We record the only property of $W_d$ needed below.  Let
$\delta_n=2^{-n}$ and let $J_n$ be the real interval
$[1-\delta_n,1-\delta_n/2]$.  Choose fixed, sufficiently small enlargements
$Q_n\Subset\mathbb D$ of the $\delta_n$-neighborhoods of $J_n$.  On $Q_n$,
\begin{equation}\label{eq:comparabilities}
 A\asymp\delta_n,\qquad b\asymp\delta_n^2.
\end{equation}
Restricting the integral in \eqref{eq:marginal} to $b\leq u\leq2b$ gives
\begin{equation}\label{eq:Wlower}
 W_d(\lambda)\geq\frac{c_d'}{\delta_n},
 \qquad \lambda\in Q_n.
\end{equation}
Indeed, on that interval the integrand is bounded below by a dimensional
constant times $\delta_n b^{-2}$, and the interval has length $b$.

For each $r\in J_n$, the submean inequality for the analytic function $h'$
on a disk of radius comparable to $\delta_n$ gives
\[
 |h'(r)|^2\leq C\delta_n^{-2}
 \int_{Q_n}|h'(\lambda)|^2dA(\lambda).
\]
Combining this with \eqref{eq:Wlower} and then integrating over $J_n$ yields
\begin{equation}\label{eq:lineenergy}
 \int_{J_n}|h'(r)|^2dr
 \leq C_d\int_{Q_n}|h'(\lambda)|^2W_d(\lambda)dA(\lambda)
 =:C_d E_n.
\end{equation}
The boxes $Q_n$ have bounded overlap.  Therefore
\begin{align*}
 \int_{1/2}^1|h'(r)|dr
 &\leq\sum_n |J_n|^{1/2}
 \left(\int_{J_n}|h'(r)|^2dr\right)^{1/2}\\
 &\leq C_d\left(\sum_n\delta_n\right)^{1/2}
 \left(\sum_n E_n\right)^{1/2}
 \leq C_d\mathcal E_{e_1}(Qf)^{1/2}.
\end{align*}
Thus $h(r)$ has a finite limit at $1$.  The same submean estimate on slightly
smaller boxes controls the oscillation in every fixed Stolz cone, so this
limit is the tip trace.  Finally, bounded point evaluation at $e_1/2$ in
$H^2(\Bd)$ gives
\[
 |\Tr_{e_1}f|
 \leq |h(1/2)|+C_d\mathcal E_{e_1}(Qf)^{1/2}
 \leq C_d\bigl(\lVert Qf\rVert_{H^2}+\mathcal E_{e_1}(Qf)^{1/2}\bigr).
\]
Use \eqref{eq:contractions} and square the estimate.
\end{proof}

\section{Full affirmative solution}

\begin{theorem}[Arbitrary Poisson measures are Carleson]\label{thm:main}
Let $d\geq2$ and let $\mu$ be any finite positive Borel measure on $\Sd$.
Then there is a dimensional constant $C_d$ such that every
$f\in\mathscr D(\mu)$ satisfies
\begin{equation}\label{eq:mainestimate}
 \int_{\Bd}|f(z)|^2P[\mu](z)dV(z)
 \leq C_d\left(
 \mu(\Sd)\lVert f\rVert_{H^2(\Bd)}^2
 +\int_{\Bd}\lVert\nabla f(z)\rVert^2P[\mu](z)dV(z)
 \right).
\end{equation}
Consequently $P[\mu]dV$ is a Carleson measure for $\mathscr D(\mu)$, every
coordinate function $z_j$ is a multiplier of $\mathscr D(\mu)$, and the
tuple of coordinate multipliers is a joint $2$-isometry.
\end{theorem}

\begin{proof}
Tonelli's theorem applied to \eqref{eq:poisson} gives
\begin{align}
 \int_{\Bd}|f|^2P[\mu]dV
 &=\int_{\Sd}\Ieta(f)d\mu(\eta),\label{eq:tonelliI}\\
 \int_{\Bd}\lVert\nabla f\rVert^2P[\mu]dV
 &=\int_{\Sd}\Eeta(f)d\mu(\eta).\label{eq:tonelliE}
\end{align}
The second quantity is finite for $f\in\mathscr D(\mu)$, so
$\Eeta(f)<\infty$ for $\mu$-almost every $\eta$.  Lemmas
\ref{lem:tip} and \ref{lem:trace}, with constants independent of $\eta$, give
\[
 \Ieta(f)\leq C_d
 \bigl(\lVert f\rVert_{H^2(\Bd)}^2+\Eeta(f)\bigr)
\]
for almost every $\eta$.  Integrating this inequality and using
\eqref{eq:tonelliI}--\eqref{eq:tonelliE} proves
\eqref{eq:mainestimate}.  In particular,
\[
 \int_{\Bd}|f|^2P[\mu]dV
 \leq C_d\max\{\mu(\Sd),d\}\lVert f\rVert_{\mathscr D(\mu)}^2.
\]
This is exactly published Proposition 3.5(ii), so its equivalence proves that
all $z_j$ are multipliers.  Published Theorem 3.4 then says that their tuple
is a joint $2$-isometry.
\end{proof}

\begin{remark}
The proof works directly for arbitrary $f\in\mathscr D(\mu)$.  It does not
pass through polynomials, radial dilations, or an unproved density statement.
The argument is scalar, matching Question 1.2(a).  An operator-valued
semispectral extension would require a separate noncommutative trace lemma
and is not claimed here.
\end{remark}

\section{Adversarial verification}

\begingroup\small

\begin{center}
\begin{tabular}{@{}p{0.24\textwidth}p{0.12\textwidth}p{0.54\textwidth}@{}}
\toprule
Step & Verdict & Check \\
\midrule
Atomic disintegration & Valid & The power $\rho^{1-2d}$ in $P$ cancels the
real $2d$-dimensional polar Jacobian exactly, giving
\eqref{eq:jacobian}. \\
Tip trace & Valid & Radial energy gives ray limits; the
$\rho^{-2}$ angular energy and the Poincare inequality on compact angular
sets force a single common value. \\
Anchored Hardy bound & Valid & The scalar kernel factors as
$A_L(s)/(L^2(L-s))=(1-t)(2t+1)/6\leq3/16$. \\
Stabilizer averaging & Valid & Haar averaging fixes the radial restriction,
preserves the atomic weight, and contracts both quadratic norms by Jensen. \\
Marginal estimate & Valid & On dyadic Stolz boxes, restricting
\eqref{eq:marginal} to $u\in[b,2b]$ gives $W_d\gtrsim\delta^{-1}$ uniformly
in scale. \\
Arbitrary functions & Valid & Tonelli is applied to the actual finite
Dirichlet energy; no polynomial approximation or boundary value with respect
to $\mu$ is assumed. \\
Passage to multipliers & Valid & The final estimate is precisely the
condition already proved equivalent to coordinate multiplication in source
Proposition 3.5. \\
\bottomrule
\end{tabular}
\end{center}

The included script \texttt{code/verify\_atomic\_geometry.py} is a guard
against algebraic and scaling mistakes, not part of the proof.  It checks 800
centered-polar identities, 10,001 samples of the exact Hardy kernel, and 567
dyadic marginal samples in complex dimensions $2$ through $8$.  From the
packet directory, run:
\begin{quote}\small\ttfamily
conda run --no-capture-output -n sandbox python code/verify\_atomic\_geometry.py
\end{quote}

\section{Novelty check, limitations, and recommendation}

\paragraph{Bounded literature search.}
The four lightweight indexes for run \texttt{fa\_banach\_001} contained no
entry for arXiv:2001.09616 or this question.  External searches used the exact
unresolved sentence, the exact title with \texttt{multiplier},
\texttt{Poisson integral}, and \texttt{Carleson measure}, the published
problem/proposition labels, and the later visible Dirichlet-type papers and
publication lists of the three source authors.  The search found the original
JFA paper, papers on standard radial Dirichlet scales, and later work on the
bidisc, but no explicit answer to the arbitrary Euclidean-Poisson-measure
question.  This was a bounded search, not an exhaustive priority claim.

\paragraph{Limitations.}
The result answers the scalar arbitrary-measure problem exactly as posed.  It
does not claim an operator-valued semispectral analogue, an invariant-Poisson
analogue, sharp constants, compactness of the embedding, or multiplier
characterizations beyond the coordinate functions.

\paragraph{Recommendation.}
The mathematical verdict is \textbf{likely valid}, with high proof confidence.
Send to human review.  The most important checks are the finite-energy tip
trace argument in Lemma \ref{lem:tip} and the identification of the
stabilizer-averaged nontangential trace in Lemma \ref{lem:trace}.  The main
bibliographic task is a citation search for later papers that cite the 2020
JFA article and address Question 1.2(a).

\begin{thebibliography}{9}
\bibitem{CGR2020}
S.~Chavan, R.~Gupta, and Md.~R.~Reza,
\emph{Dirichlet-type spaces on the unit ball and joint $2$-isometries},
J.~Funct.~Anal. 279 (2020), no.~10, 108733; arXiv:2001.09616v2.
\end{thebibliography}
\endgroup

\end{document}
